\section{Vulnérabilité et sécurité des données}
La sécurité informatique de Cujas s'est construite sur un modèle de séparation défensive. L'infrastructure physique de la bibliothèque Cujas repose sur des serveurs fermés, un SFTP nocturne dictant et contraignant l'organisation humaine du travail, forçant les agents à travailler de manière isolée sur leurs postes plutôt que de manière collaborative. En effet, la sécurité et la conservation étaient des priorités importantes dans le \enquote{deuxième âge du numérique}. Avec l'irruption du web et la numérisation massive, cela multiplie le \enquote{potentiel numérique en matière de disponibilité des corpus}. Face à cet essor, les institutions patrimoniales sont aménées à réfléchir à des moyens permettant l'ouverture de leurs données tout en les protégeant. Le \enquote{deuxième âge du numérique}, terme employé par Emmanuelle Bermès, illustre les enjeux de cette période notamment face aux \enquote{enjeux de structuration de l'information au format numérique, d'interopérabilité et de pérennisation}. Cela permet aux institutions patrimoniales de mettre à disposition leurs contenus numérisés.

De plus, pour permettre cette diffusion du savoir, il est nécessaire d'utiliser les mêmes standards de structuration de la donnée, comme le MARC, utilisé pour rédiger les notices bibliographiques. Aujourd'hui, ces règles permettent la transmission des documents numérisés entre les différentes bibliothèques numériques. Toutefois, ces concepts ont été mis en place d'une façon si stricte dans la bibliothèque Cujas qu'ils paralysent aujourd'hui le travail collaboratif des agents et figent les métadonnées. 

Cette hypervigilance de la part de la bibliothèque Cujas résulte d'un traumatisme de cyberattaque révélant ainsi une vulnérabilité locale. En 2017, la bibliothèque Cujas a subi une tentative de ransomware par cryptage. L'attaque a entraîné la \enquote{perte de fichiers, et quelques documents récupérés grâce au CINES}. Heureusement, \enquote{l'essentiel des fichiers images ont pu être sauvés (retrouvés dans des sotckages), mais le travail sur les fichiers d'une cinquantaine de livres (contrôle qualité, traitement...) ont dû être refaits}. La crainte de perdre à nouveau des données amène la bibliothèque Cujas à se montrer davantage vigilante. Suite à cette attaque, \enquote{des sauvegardes journalières ont été mises en place de l'espace de traitement}. Cet événement illustre bien la fragilité des données \enquote{de traitement} par rapport aux données \enquote{d'archivage}. Ce traumatisme a eu pour effet de pousser l'informatique a verrouiller la chaîne de production au détriment de l'ergonomie. 

Cette sécurisation a entraîné un alourdissement de la chaîne de numérisation amenant à une saturation de leur \enquote{baie de numérisation}. En effet, la bibliothèque Cujas utilise un espace de stockage local permettant de cloisonner les différentes étapes du processus de production. Chacun des espaces de stockage représente une tâche à effectuer. Cet espace de stockage est constitué de NAS, possédant un espace de stockage massif entre 40 et 50 téraoctets. Aujourd'hui, cet espace est saturé, encombré de répertoires multiples qui rendent la structure globale illisible. Cette saturation résulte notamment du blocage lié à l'impossibilité temporaire de verser les fichiers SIP au CINES. La bibliothèque se trouve contrainte de les stocker dans son espace interne, l'amenant à repenser sa chaîne de numérisation. La bibliothèque envisage d'intégrer par exemple des formats d'images plus léger, comme le JPEG ou JPEG 2 000 à l'instar des TIFF, libérant ainsi de l'espace de stockage. De nombreuses bibliothèques patrimoniales, comme la Bibliothèque Sainte-Geneviève (BSG) ou la Bibliothèque Sciences Po, ont également été confronté à cette question. Ils ont choisi de faire ce basculement afin de libérer l'espace, tout en leur permettant de réduire leur empreinte carbonne liée à leur projet de numérisation. Cela ne résoudrait que partiellement le problème, car qu'adviendrait-ils des anciens versements devront-ils rester en l'état. Pour l'instant la question demeure en suspens. 

Un autre des problèmes rencontré par la bibliothèque Cujas, vient de la chaîne de numérisation elle-même. Le travail des experts métiers de la bibliothèque est extrêmement segmenté. Afin d'alléger légèrement leurs travails, des moulinettes ont été crées, permettant une forme d'automatisation dans le travail des ouvrages numérisés. Ces moulinettes aident à combler des lacunes identifiées par le personnel qui a été mis en place par le responsable DIT les soulageant par exemple dans la production des notices numériques. Cependant, le moindre changement peut entraîner une casse des \enquote{moulinettes} car correspondant à un schéma particulier de travail des équipes. Elles sont donc difficilement adaptable ou demandent une profonde réflexion afin d'améliorer au mieux la production. 

En réaction à l'attaque de 2017, le responsable du DIT a énormément sécurisés les accès des équipes allant jusqu'à différencier les droits d'accès de chacun concernant les différentes étapes de production. Personne n'a accès à tous les paramètres. La publication des données ne peut plus se faire de façon simultanée car celles-ci passent désormais par un outil SFTP sécurisé tels que WinSCP avec un traitement par lot. Ces transferts se réalisent uniquement la nuit, impactant nettement moins le travail des équipes selon le responsable du DIT. En journée, l'espace est occupé ce qui ralentirait davantage les traitements.

Le Département de patrimoine, de conservation et de numérisation (DPCN) travaille sur la baie de numérisation, uniquement accessibles grâce aux codes d'accès sur des postes prédéfinis afin d'effectuer une tâche précise. Ayant été à l'origine comme un espace de stockage, il a fallu l'adapter pour permettre des connexions simultanémées. Toutefois, des préccautions sont à prendre face à son instabilité, comme l'interdiction d'utiliser un moteur de recherche à l'intérieur afin de ne pas fragiliser l'infrastructure, et de protéger les données traitées. Par ailleurs, il est impossible pour les équipes de s'y connecter à distance que ce soit dans un autre lieu de la bibliothèque ou chez eux. Le manque de structure suffisant de cette baie oblige à ce niveau prudence afin d'éviter la corruption des données.
 
Cette surprotection des équipes de Cujas via leur chaîne de numérisation amène à réfléchir à la notion même de sécurisation autour de la donnée. En effet, à une époque où l'on pense à ouvrir celle-ci permettant à d'autres institutions patrimoniales de pouvoir l'utiliser pour leur propre collection. Cette archive ouverte interroge sur le maintien des conditions de sécurités et de préservation nécessaire, sans complètement bloquer l'échange des métadonnées entre les institutions patrimoniales. On observe une amplification de ce phénomène avec l'arrivée du Web sémantique. Bien sûr, la sécurité et la préservation sont des conditions nécessaires et importantes pour les bibliothèques mais ne doivent pas pour autant entraîner un isolement des données de l'institution à cause d'un modèle de sécurité trop stricte. Ce modèle de sécurité utilisé par Cujas est typique des silos du \enquote{deuxième âge}, créant une contradiction avec les méthodes de travail modernes empêchant une autonomie relative de la donnée. 

Cette hypervigilance justifiée à bien des égards entraîne un isolement de la bibliothèque Cujas par rapport aux autres bibliothèques numériques. La rigidité de leur condition de travail entraîne comme conséquence l'impossibilité de faire évoluer leur système pour répondre aux attentes d'aujourd'hui. Les exigences de maintenant sont davantages flexible avec le Web de données que dans les années 1990 et 2010 empêchant la bibliothèque d'améliorer l'ergonomie de son interface.